\section{Related Work}
\label{sec:related}

\subsection{Informal and Lexical Definitions}

The word itself already hints at the structure this paper wants to make
explicit. \emph{Data} comes from the Latin \emph{datum}, ``that which is
given''---a datum is received, not derived; a starting point for
reasoning, not its product~\autocite{websters1934}. Standard references
keep that sense without ever formalizing it. Rob and Coronel call data
raw, unprocessed facts~\autocite{rob2013}; Ratzan treats
\emph{information}, not data, as the organized, coherent thing, leaving
data as whatever sits underneath it~\autocite{ratzan2004}; UNESCO comes
closest to a structural requirement when it insists that data be
``expressed in a standardized way, suitable for being communicated,
interpreted, or processed''~\autocite{unesco2008}---which is a real
constraint, but not yet one stated in terms anyone could check
mechanically.

\subsection{Floridi's General Definition of Information}

Floridi's General Definition of Information treats information as data
that happen to be well-formed, meaningful, and---in some versions of the
definition---true, and it treats data themselves as little more than a
primitive: a ``lack of uniformity'' between two states of a
system~\autocite{floridi2010}. That move is deliberate on Floridi's part;
the GDI is built to explain how information arises \emph{from} data, not
to specify what data are made of. The function-theoretic definition fills
in exactly the structural detail the GDI leaves out by design, at the
layer underneath where Floridi's own account begins.

\subsection{Zins' Faceted Classification}

Zins surveys a wide range of ways people have tried to define data,
information, and knowledge, and sorts them into facets---some treat data
as \emph{objects}, others as \emph{signals}, others as
\emph{representations} of facts~\autocite{zins2007}. The definition argued
for here sits squarely inside the representation facet, and not by
accident: something standing in for something else, addressably, is the
one facet a computing system actually has to instantiate in order to run
at all.

\subsection{The DIKW Hierarchy}

The data--information--knowledge--wisdom hierarchy is one of the more
widely cited ways of describing how raw data get progressively enriched
through organization and interpretation~\autocite{ackoff1989,rowley2007}.
What it is almost never given is a precise account of its own lowest
rung---what ``data'' means at the bottom of that hierarchy is usually
assumed rather than stated. That is the gap this paper is aimed at: a
characterization of exactly that tier, one every DIKW-style account
already presupposes whether it says so or not.